% Generated by scripts/assemble.py; corrections belong in source/corrections.json.
\examhead{79}{东北师范大学 2025 年数学分析试卷}{77}
% p078-b18; source PDF page 78
\par\Needspace{5\baselineskip}
\blockanchor{p078-b18}
\majorquestion{一}计算题. 每题6分, 共30分.\par
% p078-b19; source PDF page 78
\blockanchor{p078-b19}
\question{1}$\lim_{n\to\infty}\left(1+\frac{1}{2}+\cdots+\frac{1}{n}\right)^{\frac{1}{n}}$\par
% p078-b20; source PDF page 78
\blockanchor{p078-b20}
\question{2}$\lim_{x\to 0}\frac{\int_0^{\sin^2x}te^t\mathrm{d}t}{\sqrt{1 + x^4} - 1}.$\par
% p078-b21; source PDF page 78
\blockanchor{p078-b21}
\question{3}$\lim_{(x,y)\to (\infty ,\infty)}\frac{x + y}{x^2 - xy + y^2}.$\par
% p078-b22; source PDF page 78
\blockanchor{p078-b22}
\question{4}$\iint_{D} [x^{2} + y^{2}] \, \mathrm{d}x \, \mathrm{d}y$, 其中 $D = \{(x, y) \mid -1 \leq x \leq 1, -1 \leq y \leq 1\}$, 函数 $[\cdot]$ 是取整函数.\par
% p079-b00; source PDF page 79
\blockanchor{p079-b00}
\question{5}$\int_{-\pi}^{\pi}\left(2 + \sum_{n = 1}^{100}\frac{\cos nx}{n}\right)^{2}\mathrm{d}x.$\par
% p079-b01; source PDF page 79
\par\Needspace{5\baselineskip}
\par\Needspace{6\baselineskip}
\blockanchor{p079-b01}
\majorquestion{二}(15 分) 计算曲线积分\par
\nopagebreak[4]
% p079-b02; source PDF page 79
\blockanchor{p079-b02}
\displayformula{\int_ {C} \left(4 x ^ {3} y ^ {3} - 3 y ^ {2} + 6 y + 1\right) \mathrm{d} x + \left(3 x ^ {4} y ^ {2} - 6 x y + 10 x - 4\right) \mathrm{d} y.}
% p079-b03; source PDF page 79
\blockanchor{p079-b03}
积分路径 C 是圆弧  $x^{2} + y^{2} = 1 (x \geq 0, y \geq 0)$ ，起点是  $(1, 0)$ ，终点是  $(0, 1)$ .\par
% p079-b04; source PDF page 79
\par\Needspace{5\baselineskip}
\blockanchor{p079-b04}
\majorquestion{三}(15 分) 设 $u = e^{x + y}$, 而 $y$ 是由方程 $y - \frac{1}{2} \sin y = x$ 所确定的函数, 求 $\frac{\mathrm{d} u}{\mathrm{d} x}, \frac{\mathrm{d}^2 u}{\mathrm{d} x^2}$.\par
% p079-b05; source PDF page 79
\par\Needspace{5\baselineskip}
\blockanchor{p079-b05}
\majorquestion{四}(15 分) 设 $f(x)$ 在 $[a, b]$ 上连续, 在 $(a, b)$ 上可导, $f(a) = 0$, 当 $x \in (a, b]$ 时, $f(x) > 0$. 证明: $\frac{f'(x)}{f(x)}$ 在 $(a, b)$ 内无界.\par
% p079-b06; source PDF page 79
\par\Needspace{5\baselineskip}
\blockanchor{p079-b06}
\majorquestion{五}(15 分) 设 $a > b > 0$, 求级数 $\sum_{n=1}^{\infty} \frac{a^n + (-b)^n}{2n} (x - 1)^n$ 的和函数.\par
% p079-b07; source PDF page 79
\par\Needspace{5\baselineskip}
\par\Needspace{6\baselineskip}
\blockanchor{p079-b07}
\majorquestion{六}(15 分) 设函数 $f(x)$ 在区间 $(- \infty, +\infty)$ 内连续可导, 令\par
\nopagebreak[4]
% p079-b08; source PDF page 79
\blockanchor{p079-b08}
\displayformula{f _ {n} (x) = \frac {n}{2} \left[ f \left(x + \frac {1}{n}\right) - f \left(x - \frac {1}{n}\right) \right], n = 1, 2, 3, \dots .}
% p079-b09; source PDF page 79
\blockanchor{p079-b09}
(1) 证明: 对任意的闭区间 $[\alpha, \beta]$, 函数列 $\{f_n(x)\}$ 在 $[\alpha, \beta]$ 上一致收敛.\par
% p079-b10; source PDF page 79
\blockanchor{p079-b10}
(2) 求 $\lim_{n\to \infty}\int_0^{2024}f_n(x)\mathrm{d}x.$\par
% p079-b11; source PDF page 79
\par\Needspace{5\baselineskip}
\blockanchor{p079-b11}
\majorquestion{七}（15分）设数项级数 $\sum_{n=1}^{\infty} a_n$ 收敛, $\sum_{n=1}^{\infty} (b_{n+1} - b_n)$ 绝对收敛, 证明: $\sum_{n=1}^{\infty} a_n b_n$ 收敛.\par
% p079-b12; source PDF page 79
\par\Needspace{5\baselineskip}
\blockanchor{p079-b12}
\majorquestion{八}(15 分) 设 $a_0 \in (-\infty, +\infty)$, $a_n = \cos a_{n-1} (n = 1, 2, 3, \cdots)$, 证明: 数列 $\{a_n\}$ 收敛.\par
% p079-b13; source PDF page 79
\par\Needspace{5\baselineskip}
\blockanchor{p079-b13}
\majorquestion{九}(15 分) 设函数 $f(x)$ 在 $[0,1]$ 上单调增加, $f(0) > 0$, $f(1) < 1$. 证明: 对任意的 $\alpha > 0$; 存在 $\xi \in (0,1)$, 使得 $f(\xi) = \xi^{\alpha}$.\par
